\documentclass[10pt,a4paper]{article}

\usepackage[utf8]{inputenc}
\usepackage[T1]{fontenc}
\usepackage[brazil]{babel}
\usepackage[margin=1.35cm]{geometry}
\usepackage{amsmath}
\usepackage{booktabs}
\usepackage{array}
\usepackage{enumitem}
\usepackage{tabularx}
\usepackage{xcolor}
\usepackage{hyperref}
\hypersetup{colorlinks=true, linkcolor=black, urlcolor=blue}

\setlength{\parindent}{0pt}
\setlength{\parskip}{2.6pt}
\setlist[itemize]{leftmargin=*, topsep=2pt, itemsep=1pt}
\setlist[enumerate]{leftmargin=*, topsep=2pt, itemsep=1pt}
\newcommand{\arquivo}[1]{\texttt{#1}}
\newcommand{\codigo}[1]{\texttt{#1}}
\newcommand{\pergunta}[1]{\textbf{#1}}
\newcolumntype{Y}{>{\raggedright\arraybackslash}X}

\title{\vspace{-1.1cm}\textbf{Guia de Defesa do Trabalho TR1}\\
\large Simulador das Camadas Física e de Enlace}
\author{Gustavo Henrique Andrade Cavalcanti \and Fernando Augusto Hortencio}
\date{Teleinformática e Redes 1 - Universidade de Brasília - 2026}

\begin{document}
\maketitle
\vspace{-0.65cm}

\begin{abstract}
Este guia consolida a explicação técnica do Simulador TR1 para a apresentação oral. O sistema transmite uma mensagem de texto entre uma thread transmissora e uma thread receptora, passando pela camada de aplicação, pela camada de enlace, pela camada física e por um meio de comunicação ruidoso. A defesa deve deixar claro que o projeto não é apenas uma interface gráfica: ele implementa, manualmente, os algoritmos centrais de enquadramento, detecção de erros, correção de erros e modulação. O professor pode perguntar tanto sobre o código quanto sobre a teoria, então o objetivo deste documento é conectar cada conceito ao arquivo e à função que o implementam.
\end{abstract}

\textbf{Palavras-chave:} camada física; camada de enlace; enquadramento; CRC-32; Hamming; QPSK; 16-QAM; AWGN.

\section{Objetivo e visão geral}

O trabalho simula, de ponta a ponta, o caminho de uma mensagem em uma rede: o texto é convertido em bits, os bits são organizados em quadros, os quadros recebem redundância para lidar com erros, e a camada física transforma esses bits em amostras de sinal elétrico em Volts. Depois disso, o meio de comunicação soma ruído gaussiano ao sinal. No receptor, o processo é invertido: o sinal é demodulado, os quadros são recuperados, erros são corrigidos ou detectados e os bits voltam a formar texto.

A ideia central de defesa é: \textbf{bits não trafegam no meio como objetos abstratos}. O que passa pelo canal é uma forma de onda, sujeita a ruído, potência, amplitude, frequência e fase. O projeto torna isso visível ao mostrar bits, quadros, sinais transmitidos, sinais recebidos, relatório de erros e texto final.

O fluxo conceitual é:
\[
\text{texto} \rightarrow \text{bits} \rightarrow \text{quadros} \rightarrow \text{sinal} \rightarrow \text{canal ruidoso} \rightarrow \text{sinal recebido} \rightarrow \text{bits} \rightarrow \text{texto}.
\]

A separação entre transmissor e receptor é feita em \arquivo{simulador.py} com duas threads e filas \codigo{queue.Queue}. O canal fica entre as duas: ele recebe o sinal limpo da thread TX, aplica ruído em \arquivo{meio\_comunicacao.py} e entrega o sinal ruidoso à thread RX. Isso atende ao espírito do enunciado: o TX e o RX são entidades separadas, e o meio é o único ponto de contato.

\section{Camada de enlace: quadros, detecção e correção}

\subsection{Enquadramento}

A camada física entrega ao receptor um fluxo contínuo de bits. A camada de enlace precisa saber onde cada quadro começa e termina. Por isso o projeto implementa três técnicas de enquadramento em \arquivo{camada\_enlace.py}.

\begin{center}
\small
\begin{tabularx}{\linewidth}{p{2.9cm}Y}
\toprule
\textbf{Técnica} & \textbf{Como funciona e o que defender} \\
\midrule
Contagem de caracteres & O quadro começa com um byte que informa quantos bytes de payload vêm em seguida. É simples e tem baixo overhead, mas se o campo de contagem for corrompido, o receptor perde o alinhamento e pode interpretar os quadros seguintes de forma errada. No código: \codigo{enquadrar\_contagem} e \codigo{desenquadrar\_contagem}. \\
FLAGs com inserção de bytes & Usa FLAG \codigo{0x7E} para delimitar quadros e ESC \codigo{0x7D} para proteger bytes especiais que aparecem no payload. Se o dado contém FLAG ou ESC, o transmissor insere ESC antes dele. No receptor, uma pequena máquina de estados distingue dado literal de delimitador. No código: \codigo{enquadrar\_bytes} e \codigo{desenquadrar\_bytes}. \\
FLAGs com inserção de bits & Usa a FLAG \codigo{01111110}. Para impedir que o payload imite essa sequência, o transmissor insere um bit 0 depois de cinco bits 1 consecutivos. O receptor remove esse 0 extra. No código: \codigo{enquadrar\_bits} e \codigo{desenquadrar\_bits}. \\
\bottomrule
\end{tabularx}
\end{center}

Perguntas prováveis: por que inserir 0 depois de cinco bits 1? Porque a FLAG tem seis bits 1 consecutivos; cortando a sequência no quinto 1, a FLAG nunca aparece dentro dos dados. Por que escapar também o ESC no byte stuffing? Porque, se ESC pudesse aparecer como dado sem ser escapado, o receptor não saberia se ele é controle ou informação.

\subsection{Detecção de erros}

Detecção de erros significa anexar ao quadro uma redundância calculada sobre os dados. No receptor, a redundância é recalculada ou validada. Se não bater, houve erro. O projeto implementa três EDCs (\textit{Error Detecting Codes}).

\begin{center}
\small
\begin{tabularx}{\linewidth}{p{2.2cm}p{1.8cm}Y}
\toprule
\textbf{EDC} & \textbf{Custo} & \textbf{Explicação de defesa} \\
\midrule
Paridade par & 1 byte no projeto & A paridade faz o número total de bits 1 ficar par. Detecta qualquer quantidade ímpar de erros, mas falha quando uma quantidade par de bits é invertida. Foi armazenada como 1 byte, e não 1 bit solto, para manter alinhamento com os enquadramentos por bytes. \\
Checksum & 16 bits & Usa soma em complemento de 1 sobre palavras de 16 bits, com \textit{end-around carry}. No receptor, a soma do payload com o checksum deve fechar em \codigo{0xFFFF}. É mais forte que paridade, mas ainda pode falhar quando erros se compensam aritmeticamente. \\
CRC-32 & 32 bits & Interpreta os bits como polinômios e usa divisão em GF(2), onde soma e subtração são XOR. O projeto implementa o CRC-32 IEEE 802 bit a bit, sem \codigo{zlib}. Usa o polinômio refletido \codigo{0xEDB88320}, inicialização \codigo{0xFFFFFFFF} e XOR final. O vetor ``123456789'' gera \codigo{0xCBF43926}. \\
\bottomrule
\end{tabularx}
\end{center}

A ordem no transmissor é importante: primeiro o projeto adiciona o EDC e depois aplica Hamming. Assim, o Hamming protege também os bits do verificador. No receptor, a ordem é invertida: primeiro tenta corrigir com Hamming, depois verifica o EDC.

\subsection{Correção com Hamming(8,4)}

O projeto usa Hamming(8,4) estendido, também chamado SECDED: \textit{Single Error Correction, Double Error Detection}. Cada bloco de 4 bits de dados vira 8 bits:
\[
p_1\ p_2\ d_1\ p_4\ d_2\ d_3\ d_4\ p_0.
\]
Os bits \(p_1\), \(p_2\) e \(p_4\) formam as paridades clássicas de Hamming. A combinação das paridades quebradas forma a síndrome, que aponta a posição do erro simples. O bit \(p_0\) é a paridade geral do bloco e permite detectar erro duplo.

\begin{center}
\small
\begin{tabularx}{\linewidth}{p{2.3cm}p{2.6cm}Y}
\toprule
\textbf{Síndrome} & \textbf{Paridade geral} & \textbf{Diagnóstico} \\
\midrule
0 & correta & Não há erro detectado. \\
\(\ne 0\) & quebrada & Há erro simples em uma das posições 1 a 7. O receptor inverte o bit apontado pela síndrome. \\
0 & quebrada & O erro está apenas no bit \(p_0\). Os dados estão íntegros. \\
\(\ne 0\) & correta & Há erro duplo. O código detecta, mas não corrige. \\
\bottomrule
\end{tabularx}
\end{center}

Por que Hamming(8,4), e não Hamming(7,4)? Porque o bit extra \(p_0\) permite detectar erro duplo, além de manter alinhamento em bytes. O custo é alto: 4 bits viram 8 bits, ou seja, a redundância dobra o tamanho do payload protegido.

\section{Camada física: de bits para sinais}

\subsection{Codificação em banda-base}

Na banda-base, os bits são representados diretamente por níveis de tensão. O projeto implementa três códigos de linha em \arquivo{camada\_fisica.py}.

\begin{itemize}
\item \textbf{NRZ-Polar:} bit 1 vira \(+V\) e bit 0 vira \(-V\). A demodulação usa a média das amostras de cada bit: média positiva indica 1, média negativa ou nula indica 0.
\item \textbf{Manchester:} sempre há transição no meio do bit, o que ajuda no sincronismo. A convenção adotada é a de Tanenbaum/G.E. Thomas: 0 é transição baixo-para-alto e 1 é alto-para-baixo. A demodulação compara a média da primeira metade do bit com a da segunda metade.
\item \textbf{Bipolar AMI:} bit 0 vira 0 V; cada bit 1 alterna entre \(+V\) e \(-V\). Isso reduz a componente DC do sinal. A demodulação usa \(|\text{média}| > V/2\) para decidir se houve um 1.
\end{itemize}

\subsection{Modulação por portadora}

Na modulação por portadora, os bits controlam amplitude, frequência ou fase de uma senoide. A ideia que unifica QPSK e 16-QAM no projeto é a representação I/Q:
\[
s(t)=I\cos(2\pi ft)-Q\sin(2\pi ft).
\]
Escolher um par \((I,Q)\) é escolher um ponto da constelação. A função \codigo{onda(i,q,ciclos)} gera o símbolo correspondente. No receptor, \codigo{correlacionar} multiplica o símbolo recebido por cosseno e seno de referência e soma os produtos. Como seno e cosseno são ortogonais ao longo de ciclos inteiros, a correlação recupera as componentes I e Q.

\begin{itemize}
\item \textbf{ASK, 1 bit/símbolo:} varia a amplitude. O bit 1 transmite portadora com amplitude \(V\); o bit 0 transmite 0 V. É simples, mas sensível a ruído de amplitude.
\item \textbf{FSK, 1 bit/símbolo:} varia a frequência. O projeto usa 2 ciclos/símbolo para 0 e 4 ciclos/símbolo para 1. O receptor compara a energia nas duas frequências.
\item \textbf{QPSK, 2 bits/símbolo:} usa 4 fases e mapeamento Gray. Cada símbolo representa um par de bits. Como vizinhos diferem em 1 bit, um erro para o vizinho tende a causar apenas 1 bit errado.
\item \textbf{16-QAM, 4 bits/símbolo:} usa uma grade 4x4, variando amplitude e fase. Carrega mais bits por símbolo, mas os pontos ficam mais próximos; por isso exige melhor SNR que QPSK.
\end{itemize}

O trade-off de modulação é essencial: mais bits por símbolo aumentam eficiência espectral, mas reduzem a distância entre pontos da constelação. Ruído desloca o ponto recebido; a decisão por mínima distância funciona enquanto o deslocamento não ultrapassa a fronteira do símbolo vizinho.

\subsection{Meio e ruído}

O meio de comunicação está em \arquivo{meio\_comunicacao.py}. Ele soma \codigo{random.gauss(media, sigma)} a cada amostra do sinal. Isso modela AWGN: ruído aditivo, branco e gaussiano. A média \(x\) pode representar deslocamento DC; o desvio padrão \(\sigma\) controla a intensidade do ruído. A potência média é calculada como \(P=(1/N)\sum v^2\), assumindo carga de \(1\Omega\). Essa potência permite comparar sinal e ruído e explicar SNR.

\section{Como o código se organiza}

\begin{center}
\small
\begin{tabularx}{\linewidth}{p{3.2cm}Y}
\toprule
\textbf{Arquivo} & \textbf{O que explicar se o professor abrir o código} \\
\midrule
\arquivo{camada\_aplicacao.py} & \codigo{texto\_para\_bits} percorre bytes UTF-8 e emite bits do mais significativo ao menos significativo. \codigo{bits\_para\_texto} reagrupa bits em bytes e usa \codigo{errors="replace"} para não travar quando o ruído corrompe uma sequência UTF-8. \\
\arquivo{camada\_enlace.py} & Contém conversões bits/bytes, três enquadramentos, três EDCs, Hamming(8,4) e as fachadas \codigo{transmitir} e \codigo{receber}. No TX: divide em blocos, adiciona EDC, aplica Hamming e enquadra. No RX: desenquadra, corrige, verifica EDC e devolve bits da aplicação. \\
\arquivo{camada\_fisica.py} & Define \(V=1\), 100 amostras por bit/símbolo, códigos de linha, portadoras, constelações e demoduladores. As funções \codigo{onda}, \codigo{correlacionar}, \codigo{bits\_do\_ponto\_mais\_proximo} e \codigo{decidir\_nivel\_gray} são o núcleo das modulações por portadora. \\
\arquivo{meio\_comunicacao.py} & Soma ruído gaussiano ao sinal e calcula potência média. É onde o canal deixa de ser ideal. \\
\arquivo{simulador.py} & Cria threads TX/RX, usa filas para representar o meio, aplica ruído entre as threads e armazena resultados intermediários para a interface. \\
\arquivo{interface\_gui.py} & Implementa a GUI em GTK 3, com fallback em Tk. Permite escolher enquadramento, EDC, Hamming, modulação, ruído e tamanho de quadro, além de mostrar bits, sinais, texto recuperado e relatório por quadro. \\
\arquivo{testes.py} & Valida aplicação, enquadramentos, EDCs, CRC conhecido, Hamming com erro simples/duplo, modulações com ruído e simulações completas. \\
\bottomrule
\end{tabularx}
\end{center}

\section{Como demonstrar em sala}

\textbf{1. Comece com a tese.} ``O simulador mostra a transmissão completa: texto vira bits, bits viram quadros, quadros viram sinal, o canal injeta ruído e o receptor tenta reconstruir a mensagem.''

\textbf{2. Rode os testes.} No diretório do projeto:
\begin{verbatim}
python3 testes.py
\end{verbatim}
Explique que os testes cobrem ida e volta UTF-8, casos críticos de stuffing, paridade/checksum/CRC, CRC-32 de referência \codigo{0xCBF43926}, Hamming corrigindo 1 erro e detectando erro duplo, todas as modulações e 15 simulações completas.

\textbf{3. Abra a GUI.}
\begin{verbatim}
python3 simulador.py
\end{verbatim}
Use uma configuração segura para demonstração: enquadramento por bits, detecção CRC, correção Hamming, NRZ-Polar, QPSK, média 0 e \(\sigma=0{,}10\). Mostre a aba TX, a aba RX, os bits de enlace, o sinal transmitido, o sinal recebido e o relatório por quadro.

\textbf{4. Aumente o ruído aos poucos.} Primeiro o Hamming corrige alguns erros. Depois aparecem erros duplos ou falhas de EDC. Por fim, o texto sai com diferenças. Essa é a melhor demonstração visual do limite entre detecção, correção e capacidade do canal.

\textbf{5. Mostre três pontos do código.} Em \arquivo{camada\_enlace.py}, mostre \codigo{transmitir} e \codigo{receber}; em \arquivo{camada\_fisica.py}, mostre \codigo{onda} e \codigo{correlacionar}; em \arquivo{simulador.py}, mostre \codigo{executar\_simulacao}.

\section{Perguntas prováveis e respostas curtas}

\begin{enumerate}
\item \pergunta{O que o trabalho simula?} Uma transmissão completa entre TX e RX, passando por aplicação, enlace, física e meio ruidoso.
\item \pergunta{Por que usar threads?} Para separar transmissor e receptor. As filas representam o meio entre eles.
\item \pergunta{Onde o ruído entra?} Depois do TX e antes do RX, somado a cada amostra em \arquivo{meio\_comunicacao.py}.
\item \pergunta{Qual a ordem da camada de enlace no TX?} Divide em blocos, adiciona EDC, aplica Hamming e enquadra.
\item \pergunta{E no RX?} Desenquadra, corrige com Hamming, verifica EDC e devolve os bits da aplicação.
\item \pergunta{Por que Hamming depois do EDC?} Para proteger também os bits do verificador.
\item \pergunta{Por que Hamming(8,4)?} Porque corrige erro simples, detecta erro duplo e mantém alinhamento em bytes.
\item \pergunta{O que o CRC faz?} Divide a mensagem como polinômio em GF(2); resto diferente indica erro.
\item \pergunta{O que é GF(2)?} Aritmética módulo 2, na qual soma e subtração são XOR.
\item \pergunta{Por que paridade é fraca?} Porque erros em quantidade par podem passar despercebidos.
\item \pergunta{Por que cinco 1s no bit stuffing?} Porque a FLAG tem seis 1s; inserir 0 após cinco impede a imitação da FLAG.
\item \pergunta{O que é I/Q?} São componentes ortogonais que representam amplitude e fase de um símbolo.
\item \pergunta{Como o RX recupera I/Q?} Por correlação com cosseno e seno de referência.
\item \pergunta{Por que Gray em QPSK e 16-QAM?} Para que símbolos vizinhos difiram em apenas 1 bit.
\item \pergunta{Qual modulação é mais sensível a ruído?} ASK é muito sensível a amplitude; 16-QAM também exige SNR maior por ter pontos mais próximos.
\item \pergunta{Por que 16-QAM carrega 4 bits por símbolo?} Porque possui 16 pontos e \(\log_2(16)=4\).
\item \pergunta{O que é SNR no projeto?} Relação entre potência média do sinal e potência média do ruído.
\item \pergunta{O que provar se perguntarem sobre corretude?} Rodar \codigo{python3 testes.py} e citar os testes de CRC, Hamming, stuffing e simulação completa.
\end{enumerate}

{\footnotesize \textbf{Fontes locais usadas:} código e estudos do repositório \arquivo{TrabalhoTR1}; materiais de camada física e enlace do repositório de Teleinformática e Redes; relatório LaTeX original do grupo; suíte \arquivo{testes.py}.}

\end{document}
